\section*{\centerline{Foreword}}

It is with great pleasure that I introduce this thesis, which presents a significant contribution into one of the most interesting mysteries of physics: The nature of dark matter through accelerator based research.
Conducted within the framework of the Compact Muon Solenoid (CMS) experiment at the Large Hadron Collider (LHC) at CERN, this work is an example of the cutting-edge science being pursued in the mission 
to understand the fundamental constituents of our universe.
The author of this thesis, Dr. Alp Akpinar, has undertaken the challenging task of searching 
and constraining hypothesized interactions between the Higgs boson and dark matter. 
These measurements are among the most precise to date and provide some of the most 
stringent constraints, pushing the boundaries of our current knowledge and 
setting the stage for future explorations and discoveries.

Beyond the immediate findings, this thesis also takes an innovative approach by 
exploring the future of dark matter analyses in CMS. 
The incorporation of deep learning models into these analyses represents a 
forward-thinking strategy to enhance sensitivity and improve the potential 
for discovery in the next generation searches. 
Furthermore, the design and development of a new data processing hardware for the 
Inner Tracker of CMS, planned for deployment in the next phase of the 
LHC (so-called High Lumi LHC), highlight the Dr. Akpinar's commitment 
to advancing detector technologies.

I extend my heartfelt congratulations to Dr. Alp Akpinar for his outstanding work 
and wish him every success in their future endeavors.

\vspace{20pt}

Dr. Zeynep Demiragli

Associate Professor of Physics, Boston University